\documentclass[12pt]{article}
\usepackage{amsmath,amssymb,amsfonts}
\usepackage[margin=1in]{geometry}
\begin{document}
\begin{center}
{\large\bfseries 28th Irish Mathematical Olympiad\\25 April 2015}
\end{center}
\bigskip\noindent\textbf{Paper 1}\bigskip
\begin{enumerate}
\item In the triangle $ABC$, the length of the altitude from $A$ to $BC$ is equal to 1. $D$ is the midpoint of $AC$. What are the possible lengths of $BD$?

\item A regular polygon with $n \geq 3$ sides is given. Each vertex is coloured either red, green or blue, and no two adjacent vertices of the polygon are the same colour. There is at least one vertex of each colour.

Prove that it is possible to draw certain diagonals of the polygon in such a way that they intersect only at the vertices of the polygon and they divide the polygon into triangles so that each such triangle has vertices of three different colours.

\item Find all positive integers $n$ for which both $837 + n$ and $837 - n$ are cubes of positive integers.

\item Two circles $\mathcal{C}_1$ and $\mathcal{C}_2$, with centres $D$ and $E$ respectively, touch at $B$. The circle having $DE$ as diameter intersects the circle $\mathcal{C}_1$ at $H$ and the circle $\mathcal{C}_2$ at $K$. The points $H$ and $K$ both lie on the same side of the line $DE$. $HK$ extended in both directions meets the circle $\mathcal{C}_1$ at $L$ and meets the circle $\mathcal{C}_2$ at $M$. Prove that
\begin{enumerate}
\item[(a)] $|LH| = |KM|$;
\item[(b)] the line through $B$ perpendicular to $DE$ bisects $HK$.
\end{enumerate}

\item Suppose a doubly infinite sequence of real numbers
\[
\ldots,\, a_{-2},\, a_{-1},\, a_0,\, a_1,\, a_2,\, \ldots
\]
has the property that
\[
a_{n+3} \;=\; \frac{a_n + a_{n+1} + a_{n+2}}{3}, \qquad \text{for all integers } n.
\]
Show that if this sequence is bounded (i.e., if there exists a number $R$ such that $|a_n| \leq R$ for all $n$), then $a_n$ has the same value for all $n$.
\end{enumerate}

\newpage
\begin{center}
{\large\bfseries 28th Irish Mathematical Olympiad\\25 April 2015}
\end{center}
\bigskip\noindent\textbf{Paper 2}\bigskip
\begin{enumerate}
\setcounter{enumi}{5}
\item Suppose $x$, $y$ are nonnegative real numbers such that $x + y \leq 1$. Prove that
\[
8xy \leq 5x(1-x) + 5y(1-y),
\]
and determine the cases of equality.

\item Let $n > 1$ be an integer and $\Omega := \{1, 2, \ldots, 2n-1, 2n\}$ the set of all positive integers that are not larger than $2n$.

A nonempty subset $S$ of $\Omega$ is called \emph{sum-free} if, for all elements $x$, $y$ belonging to $S$, $x + y$ does not belong to $S$. We allow $x = y$ in this condition.

Prove that $\Omega$ has more than $2^n$ distinct sum-free subsets.

\item In triangle $\triangle ABC$, the angle $\angle BAC$ is less than $90^\circ$. The perpendiculars from $C$ on $AB$ and from $B$ on $AC$ intersect the circumcircle of $\triangle ABC$ again at $D$ and $E$ respectively. If $|DE| = |BC|$, find the measure of the angle $\angle BAC$.

\item Let $p(x)$ and $q(x)$ be non-constant polynomial functions with integer coefficients. It is known that the polynomial
\[
p(x)q(x) - 2015
\]
has at least 33 different integer roots. Prove that neither $p(x)$ nor $q(x)$ can be a polynomial of degree less than three.

\item Prove that, for all pairs of nonnegative integers $j$, $n$,
\[
\sum_{k=0}^{n} k^j \binom{n}{k} \;\geq\; 2^{n-j} n^j.
\]
\end{enumerate}
\end{document}
